% =============================================================================
% Section 7b: Compensability Analysis
% =============================================================================
\section{Compensability Analysis}\label{sec:compensability}

\subsection{Motivation}

The ISI composite is defined as an arithmetic mean of six axis scores
(\cref{eq:composite}).  Arithmetic-mean aggregation is fully
\emph{compensatory}: a high score on one axis can offset a low score
on another, and both profiles produce the same composite.  This
property has important policy implications.  A country with a
moderate ISI may be moderately concentrated on all axes (genuinely
moderate risk), or it may be extremely concentrated on one axis and
diversified on all others (a hidden single-axis vulnerability masked
by the composite).

This section quantifies the degree of internal compensation for each
country using two metrics:

\begin{enumerate}[leftmargin=2em]
  \item \textbf{Max--composite gap.}  Define $G_i = \max_j A_{ij} - \text{ISI}_i$.
    A large gap indicates that the country's highest-axis score is
    far above its composite, meaning that the mean aggregation is
    substantially compressing a peak vulnerability.
  \item \textbf{Coefficient of variation (CV).}  The standard deviation
    of the six axis scores divided by their mean:
    $\text{CV}_i = \sigma_i / \bar{A}_i$.  A high CV indicates that
    the country's axis profile is highly uneven---i.e., the degree of
    internal compensation is large.
\end{enumerate}

\subsection{Results}

% Table T17: Compensability analysis
\begin{longtable}{@{}rl S[table-format=1.8] S[table-format=1.8] S[table-format=1.8] S[table-format=1.8] S[table-format=1.4]@{}}
\caption{Compensability analysis.  Gap $=$ Max $-$ Composite.  CV: coefficient of
variation of the six axis scores per country.  High Gap and CV indicate strong
internal compensation where low axes offset high axes.}
\label{tab:compensability}\\
\toprule
{\textbf{Rk}} & {\textbf{Country}} & {\textbf{Max}} & {\textbf{Min}} & {\textbf{Range}} & {\textbf{Gap}} & {\textbf{CV}} \\
\midrule
\endfirsthead
\caption[]{Compensability (continued).}\\
\toprule
{\textbf{Rk}} & {\textbf{Country}} & {\textbf{Max}} & {\textbf{Min}} & {\textbf{Range}} & {\textbf{Gap}} & {\textbf{CV}} \\
\midrule
\endhead
\bottomrule
\endlastfoot
  1 & Malta & 1.00000000 & 0.13998786 & 0.86001214 & 0.48251852 & 0.6810 \\
  2 & Cyprus & 0.94106062 & 0.12407171 & 0.81698891 & 0.47284798 & 0.6144 \\
  3 & Ireland & 0.61373023 & 0.14666620 & 0.46706403 & 0.18570504 & 0.3759 \\
  4 & Denmark & 0.72368274 & 0.12401663 & 0.59966611 & 0.29925536 & 0.5311 \\
  5 & Croatia & 0.74648317 & 0.17343982 & 0.57304335 & 0.34213549 & 0.4592 \\
  6 & Finland & 0.79733123 & 0.11775073 & 0.67958050 & 0.39527987 & 0.5575 \\
  7 & Sweden & 0.88107943 & 0.11646504 & 0.76461439 & 0.48236411 & 0.7695 \\
  8 & Austria & 0.58506379 & 0.14903467 & 0.43602912 & 0.20256573 & 0.4000 \\
  9 & Italy & 0.92057651 & 0.14103384 & 0.77954267 & 0.54671565 & 0.7164 \\
  10 & Slovenia & 0.66661871 & 0.11917574 & 0.54744297 & 0.29682685 & 0.4682 \\
  11 & Greece & 0.71982808 & 0.16272142 & 0.55710666 & 0.35445458 & 0.5585 \\
  12 & Luxembourg & 0.64483831 & 0.11462818 & 0.53021013 & 0.28466938 & 0.5696 \\
  13 & Estonia & 0.46987927 & 0.18172493 & 0.28815434 & 0.12655750 & 0.3074 \\
  14 & Bulgaria & 0.86064518 & 0.14843767 & 0.71220751 & 0.52036534 & 0.7358 \\
  15 & Netherlands & 0.95955396 & 0.12107826 & 0.83847570 & 0.62439694 & 0.8804 \\
  16 & Hungary & 0.54471981 & 0.12222969 & 0.42249012 & 0.21505754 & 0.4482 \\
  17 & Portugal & 0.50778629 & 0.18303606 & 0.32475023 & 0.18605199 & 0.3729 \\
  18 & Romania & 0.56353195 & 0.15814280 & 0.40538915 & 0.24489003 & 0.4224 \\
  19 & Czechia & 0.52948100 & 0.16116762 & 0.36831338 & 0.21151470 & 0.4825 \\
  20 & Poland & 0.48860717 & 0.13314812 & 0.35545905 & 0.19299396 & 0.4613 \\
  21 & Belgium & 0.44734854 & 0.15348744 & 0.29386110 & 0.15938073 & 0.3512 \\
  22 & Germany & 0.58266751 & 0.09762824 & 0.48503927 & 0.31478057 & 0.6123 \\
  23 & Lithuania & 0.40227798 & 0.13056067 & 0.27171731 & 0.13687278 & 0.3973 \\
  24 & Spain & 0.43235135 & 0.11235234 & 0.31999901 & 0.17141158 & 0.4575 \\
  25 & Latvia & 0.36822299 & 0.12316988 & 0.24505311 & 0.12127162 & 0.4526 \\
  26 & Slovakia & 0.48550648 & 0.00000000 & 0.48550648 & 0.24909081 & 0.6849 \\
  27 & France & 0.37034617 & 0.10039135 & 0.26995482 & 0.13467491 & 0.4728 \\
\end{longtable}

\Cref{tab:compensability} reports both metrics for all 27~countries,
ordered by composite rank.  The five countries with the largest
max--composite gaps are:

\begin{enumerate}[leftmargin=2em]
  \item \textbf{Netherlands} (gap = 0.624, CV = 0.880).  The
    Netherlands has the most internally compensated profile in the
    dataset: its defence score (0.960) is nearly three times its
    composite (0.335).  The arithmetic mean conceals a critical
    single-axis vulnerability in defence procurement concentration.
  \item \textbf{Italy} (gap = 0.547, CV = 0.716).  Italy's defence
    axis (0.921) dominates its profile, while financial (0.141),
    technology (0.151), and critical inputs (0.183) are all near the
    EU-27 minimums.
  \item \textbf{Bulgaria} (gap = 0.520, CV = 0.736).  Defence (0.861)
    is by far the highest axis; financial (0.176) and technology
    (0.170) are below the EU-27 mean.
  \item \textbf{Malta} (gap = 0.483, CV = 0.681).  Despite ranking
    first overall, Malta's gap is large because its defence and
    logistics scores (both 1.000) are far above its financial (0.140)
    and technology (0.205) scores.
  \item \textbf{Sweden} (gap = 0.482, CV = 0.770).  Sweden's
    logistics axis (0.881) is the primary driver, with financial
    (0.116) and technology (0.145) well below the composite.
\end{enumerate}

\subsection{Interpretation}

The compensability analysis reveals a systematic pattern: countries
with the highest max--composite gaps are overwhelmingly those whose
profiles are dominated by the defence or logistics axes.  Of the ten
countries with the largest gaps, nine have either defence or logistics
as their maximum axis.

This has two implications for index interpretation:

\begin{enumerate}[leftmargin=2em]
  \item \textbf{The composite understates peak vulnerability.}
    For countries like the Netherlands (gap = 0.624), the composite
    ISI of 0.335 conveys ``moderate concentration,'' but the peak
    defence-axis score of 0.960 conveys ``near-maximum concentration
    on a single critical dimension.''  Policy-makers should examine
    axis-level profiles alongside the composite to avoid false
    reassurance.

  \item \textbf{CV discriminates policy archetypes.}  Countries can
    be classified by their CV into ``balanced'' profiles
    (CV $< 0.40$: Estonia, Ireland, Austria) and ``peaked'' profiles
    (CV $> 0.60$: Netherlands, Sweden, Italy, Bulgaria, Malta,
    Cyprus).  Balanced-profile countries face diffuse risk across
    multiple axes; peaked-profile countries face concentrated risk
    on one or two axes.  The policy prescriptions differ accordingly:
    balanced countries need broad-based diversification, while peaked
    countries need targeted axis-specific interventions.
\end{enumerate}

\subsection{Relationship to rank volatility}

There is a strong positive association between compensability and
rank volatility under Dirichlet weight perturbation
(\cref{sec:dirichlet}).  Countries with the highest CV values
(Netherlands: 0.880, Sweden: 0.770, Bulgaria: 0.736) are also among
those with the highest rank standard deviations under random
weighting (Netherlands: 3.07, Bulgaria: 2.03).  This is intuitive:
a peaked profile means that the country's composite depends heavily
on the weight assigned to its dominant axis, making the ranking
sensitive to any weighting scheme that departs from equal weights.

Conversely, countries with low CV values (Estonia: 0.307, Ireland:
0.376, Austria: 0.400) exhibit low rank volatility
($\sigma_{\text{rank}} < 1.5$), because their relatively balanced
profiles produce similar composites under any reasonable weighting.

This relationship between compensability and rank volatility provides
independent validation of both the Dirichlet Monte Carlo exercise
and the compensability metrics: each captures the same underlying
structural feature---profile unevenness---through a different lens.
